\documentclass[12pt]{article}
\usepackage[utf8]{inputenc}
\usepackage[T1]{fontenc}
\usepackage{hyperref}
\title{Ukue Language Resources for K-12 Curriculum Development}
\author{Compiled by Claude Sonnet 5 (claude-sonnet-5) for Andr\'as Kornai}
\date{}
\begin{document}
\maketitle

\section{Introduction}
Ukue (also called Epinmi; ISO 639-3 uku; Glottocode ukue1238) is an Edoid language of Nigeria, with about 14,000 speakers reported for 2000. Wikipedia places it in Ondo State; Joshua Project places the related Ukue-Ehuen people mainly in Ekiti State. Some sources treat Ukue as the same language as Ehuen. It has seven vowels in two harmonic sets and fricatives but no sibilants. Search turned up only general reference material.

\section{Primary Educational Resources}
No primers, textbooks or school programmes for Ukue were found in search results.

\section{Linguistic and Reference Materials}
\begin{itemize}
\item Wikipedia, ``Ukue language'': \url{https://en.wikipedia.org/wiki/Ukue_language}
\item Wikipedia mirror, ISO 639:uku: \url{https://wikipedia.com/wiki/ISO_639:uku}
\item Wikipedia, ``Edoid languages'': \url{https://en.wikipedia.org/wiki/Edoid_languages}
\item Journal article hosted at unizik.edu.ng (title not shown in results; relevance unverified): \url{https://journals.unizik.edu.ng/ajill/article/download/153/132/289}
\end{itemize}

\section{Religious and Community Resources}
\begin{itemize}
\item Joshua Project, Ukue-Ehuen in Nigeria: \url{https://joshuaproject.net/people_groups/15697}
\end{itemize}

\section{Digital and Audio Resources}
None found.

\section{Curriculum Development Context}
The language is not known to be taught in schools. Basic questions (relation to Ehuen, location, orthography) would need to be settled with the community before curriculum work.

\end{document}